% ============================================================
%  Analysis Section — Active/Passive Sentence Processing
%  English and Punjabi Eye-Tracking Study
%  Compile with pdflatex or xelatex
% ============================================================

\documentclass[12pt, a4paper]{article}

\usepackage[margin=2.5cm]{geometry}
\usepackage{amsmath}
\usepackage{booktabs}
\usepackage{array}
\usepackage{longtable}
\usepackage{multirow}
\usepackage{graphicx}
\usepackage{setspace}
\usepackage{natbib}
\usepackage{hyperref}
\usepackage{caption}
\usepackage{subcaption}
\usepackage{xcolor}
\usepackage{siunitx}
\usepackage{threeparttable}

\onehalfspacing

% ---- convenience macros for statistics ----
\newcommand{\pval}[1]{$p = #1$}
\newcommand{\pless}[1]{$p < #1$}
\newcommand{\tstat}[2]{$t = #1$, $p = #2$}
\newcommand{\bstat}[3]{$\hat{\beta} = #1$, $SE = #2$, $t = #3$}

\title{Active and Passive Sentence Processing in English and Punjabi:\\
       An Eye-Tracking Investigation of the Agent-First Principle}

\author{}
\date{}

\begin{document}

\maketitle

% ============================================================
\section{Method}
% ============================================================

\subsection{Participants}

Fifty participants completed the experiment (25 English-dominant, 25 Punjabi-dominant;
details reported in the Participants section). All had normal or corrected-to-normal vision.

% ============================================================
\subsection{Design and Stimuli}
% ============================================================

A 2 (Language: English, Punjabi) $\times$ 3 (Sentence Type: active, passive,
passive-with-dropped-agent) within-participants design was used.
Each participant contributed 150 trials per condition in the sentence-audio
block and an equal number of matched free-viewing trials, yielding 900
sentence-audio trials and 900 free-viewing trials in total.
The final analysis dataset was perfectly balanced: 150 trials per
Language $\times$ Sentence-Type cell (Table~\ref{tab:design}).

\begin{table}[h]
\centering
\caption{Trial counts per condition in the sentence-audio block.}
\label{tab:design}
\begin{tabular}{lll r}
\toprule
Language & Sentence Type & Agent-First Alignment & $N$ \\
\midrule
English  & Active                & Agent-first     & 150 \\
English  & Passive               & Non-agent-first & 150 \\
English  & Passive (agent dropped) & Non-agent-first & 150 \\
Punjabi  & Active                & Agent-first     & 150 \\
Punjabi  & Passive               & Non-agent-first & 150 \\
Punjabi  & Passive (agent dropped) & Non-agent-first & 150 \\
\bottomrule
\end{tabular}
\end{table}

Each scene stimulus depicted a transitive event (e.g., a woman pushing a man)
and was paired with either an active sentence (\textit{The woman is pushing the man}),
a passive sentence with an overt by-phrase (\textit{The man is being pushed by the woman}),
or an agentless passive (\textit{The man is being pushed}).
All stimuli were matched for image content across conditions.
Two regions of interest (ROIs) were defined per image: an \textit{agent} ROI
(the entity performing the action; AOI label prefix \texttt{s\_}) and an
\textit{object/patient} ROI (the entity receiving the action; prefix \texttt{o\_}).

% ============================================================
\subsection{Apparatus and Procedure}
% ============================================================

Eye movements were recorded using an EyeLink system sampling at [X] Hz.
Each trial began with a central fixation cross.
In free-viewing trials participants inspected the image in silence.
In sentence-audio trials the same image was presented concurrently with an
auditory sentence.
The free-viewing block was used to establish participant-specific and
image-specific baselines for all gaze outcomes.

% ============================================================
\section{Data Processing}
% ============================================================

\subsection{Raw Data Extraction and Standardisation}

Raw data consisted of three EyeLink Data Viewer report types exported per
participant: fixation reports, interest-area (ROI) reports, and saccade reports,
giving 50 files per report type ($N = 150$ files in total).
Files were exported in EyeLink's tab-delimited XLS format (UTF-16 encoded) and
processed using a custom Python pipeline (\texttt{process\_eye\_tracking\_reports.py}).

For each file, column names were normalised to snake\_case, and
missing-value markers used by EyeLink (\texttt{.}, \texttt{UNDEFINED},
empty string) were converted to \texttt{NaN}.
Sentence type was inferred from stimulus identifiers using regular expressions:
strings containing \textit{ACT} or \textit{active} were coded active;
\textit{PNA}, \textit{noagent}, \textit{woa}, or \textit{dropped} were coded
passive-with-dropped-agent; \textit{PWA} or \textit{passive} were coded passive.
Session type (free-viewing vs.\ sentence-audio) was inferred from a session
variable column using matching on \textit{free} and \textit{experiment}/\textit{sentence}/\textit{audio}.

Participant identifiers were extracted from filenames following the pattern
\texttt{edf[NN]}.
Trial-level features were computed by aggregating fixation, interest-area,
and saccade rows within each participant $\times$ trial cell.

\subsection{Centre-of-Screen Bias: Identification and Correction}

Inspection of the raw fixation data revealed a strong centre-of-screen
bias at the first fixation of each trial.
Specifically, 87.2\% of fixations with ordinal position index 1
(i.e., the first registered fixation per trial) fell within 100 pixels
of screen centre ($x = 960$, $y = 540$; mean first-fixation position:
$\bar{x} = 949$~px, $SD = 118$~px; $\bar{y} = 527$~px, $SD = 84$~px).
This clustering is attributable to the pre-trial central fixation cross,
which resets gaze to the screen centre before every stimulus onset.
Crucially, this fixation carries no sentence-processing signal: it reflects
a motor return to the calibration point, not a stimulus-driven oculomotor response.

To prevent this artefact from contaminating trial-level summaries,
fixation index~1 was \emph{excluded} from all trial-level aggregations
(mean fixation duration, total fixation duration, fixation count).
The raw fixation records were retained in their entirety for subsequent
temporal analyses, which use ordinal fixation position as the predictor
and are therefore explicitly indexed from position~2 onward.
Note that the interest-area (ROI dwell time) and saccade reports are
unaffected because EyeLink generates these from a separate data stream
that does not include the pre-stimulus fixation cross event.

\subsection{Baseline Subtraction}

To isolate the cognitive processing cost attributable to sentence structure,
all trial-level fixation and saccade outcomes were expressed as deviations
from free-viewing baselines using the same image.
For each participant and image, mean outcome values from the free-viewing
block were subtracted from the corresponding sentence-audio trial values,
yielding \emph{delta} ($\Delta$) outcomes.
Participant-specific baselines were available for all 900 sentence-audio
trials (100\% participant $\times$ image baseline coverage), eliminating
the need to fall back on image-level averages.
All $\Delta$ outcomes are therefore interpretable as the additional gaze
demand imposed by concurrent sentence processing relative to passive
image viewing.

\subsection{AOI Role Coding}

AOI labels were assigned roles automatically from EyeLink label prefixes:
labels beginning with \texttt{s\_} were coded as \emph{agent/subject} AOIs;
labels beginning with \texttt{o\_} were coded as \emph{object/patient} AOIs;
all remaining labels were coded as background.
Zero of 3,636 interest-area observations had a missing role assignment.

\subsection{Agent-First Coding}

Three theoretically motivated binary contrast variables were derived from
sentence type:
\begin{itemize}
  \item \textbf{agent\_first\_alignment}: \textit{agent-first} for active;
        \textit{non-agent-first} for both passive conditions.
  \item \textbf{agent\_overt}: \textit{agent-overt} for active and passive
        (by-phrase present); \textit{agent-dropped} for passive-with-dropped-agent.
  \item \textbf{passive\_cost\_contrast}: 0 for active, 1 for passive and
        passive-with-dropped-agent (tests overall non-agent-first cost).
\end{itemize}

\subsection{First Non-Centre AOI Visit}

To operationalise initial attentional orientation, the first fixation at
ordinal position $\geq 2$ that landed within a coded agent or object AOI was
extracted per trial.
The binary variable \textit{agent\_first\_look} was set to 1 if this
fixation landed in the agent AOI and 0 if it landed in the object AOI.
Trials on which no post-centre fixation landed in either AOI were excluded
from this analysis ($n = 7$ trials, $< 1\%$; final $N = 893$ trials).

\subsection{Final Analysis Datasets}

The processing pipeline produced the following datasets used in all analyses:

\begin{enumerate}
  \item \textbf{Trial-level baseline-adjusted dataset}
        (\texttt{sentence\_audio\_baseline\_adjusted.csv}):
        900 rows $\times$ 77 columns; one row per sentence-audio trial with
        $\Delta$ outcomes and agent-first coding variables.
  \item \textbf{Interest-area dataset}
        (\texttt{interest\_area\_agent\_first\_dataset.csv}):
        3,636 rows; one row per trial $\times$ AOI role, sentence-audio trials only.
  \item \textbf{First AOI visit dataset}
        (\texttt{first\_aoi\_visit\_dataset.csv}):
        893 rows; one row per trial coding the first non-centre AOI fixation.
  \item \textbf{Temporal gaze dynamics tables}
        (\texttt{prop\_agent\_by\_fixation\_index.csv}):
        Proportion of agent-directed fixations at each ordinal fixation position,
        by condition.
\end{enumerate}

% ============================================================
\section{Statistical Analysis}
% ============================================================

\subsection{Model Structure}

All primary analyses used Generalised Additive Mixed Models (GAMMs) fitted
with the \texttt{bam()} function in the \texttt{mgcv} package for R \citep{Wood2017}.
Trial-level continuous outcomes were modelled with Gaussian family; the binary
agent-first look outcome was modelled with binomial family (logit link).

The fixed-effects structure for all trial-level models was:
\begin{equation}
  \text{outcome} \sim \text{Language} \times \text{SentenceType}
    + s(\text{ParticipantID}, \, \text{bs=``re''})
    + s(\text{ImageID}, \, \text{bs=``re''})
\end{equation}
where $s(\cdot, \text{bs=``re''})$ denotes a random-effect smooth
(equivalent to a standard random intercept).
Interest-area models included an additional random smooth for individual
AOI labels, $s(\text{AOILabel}, \, \text{bs=``re''})$, to account for the
large variance attributable to specific image content.
The reference levels were English and active.
All models were fitted with fast REML (\texttt{method = "fREML"}).

Temporal gaze dynamics were modelled using a GAM with condition-specific
smooths over fixation ordinal position:
\begin{equation}
  \hat{p}(\text{agent look}) \sim
    s(\text{FixationIndex}, \, \text{by=Condition}, \, k = 8)
    + \text{Condition}
\end{equation}
weighted by the number of contributing fixations per bin.

Non-parametric screening tests (Kruskal--Wallis omnibus followed by
Mann--Whitney pairwise tests with Benjamini--Hochberg correction) were
used as pre-GAMM exploratory checks; these results are reported for
completeness but are not the primary inferential basis.

\subsection{Hypotheses}

The study addressed four pre-specified hypotheses derived from the
agent-first principle \citep[e.g.,][]{MacWhinney1982, Gernsbacher1990}:

\begin{description}
  \item[H1 — Passive processing cost.]
        Passive sentences (non-agent-first) impose a greater gaze-based
        processing cost than active sentences (agent-first), reflected in
        longer mean fixation durations, fewer fixations per trial, lower
        saccade rates, and increased dwell time in interest areas.
  \item[H2 — Agent-drop cost.]
        Passive sentences without an overt agent (passive-dropped-agent)
        impose at least as great a processing cost as passive sentences
        with an overt by-phrase, because the absence of an explicit agent
        referent further disrupts agent-first event parsing.
  \item[H3 — Cross-linguistic modulation.]
        The processing cost associated with non-agent-first structure
        differs between English and Punjabi, reflecting structural or
        morphological differences in how each language marks thematic roles.
        Specifically, Punjabi's richer morphological case marking is
        predicted to reduce the passive cost relative to English.
  \item[H4 — Agent-first gaze bias.]
        During active sentences, listeners orient gaze to the agent AOI
        earlier and for longer than during passive sentences,
        consistent with rapid integration of agent-first structural cues
        into the developing event representation.
\end{description}

% ============================================================
\section{Results}
% ============================================================

% ============================================================
\subsection{Trial-Level Processing Cost (H1, H2, H3)}
% ============================================================

\subsubsection{Mean Fixation Duration}

Passive sentence processing produced significantly longer mean fixation
durations than active processing after controlling for free-viewing baselines
($\Delta$ passive English: $M = +94.9$~ms, $SD = 474.9$; $\Delta$ active
English: $M = -52.9$~ms, $SD = 511.3$).
The GAMM confirmed this contrast:
$\hat{\beta} = 147.6$~ms, $SE = 54.3$, $t = 2.72$, $p = .007$.
Passive-with-dropped-agent also showed numerically elevated fixation
duration relative to active ($\Delta M = +5.8$~ms in English), but this
contrast was not significant ($\hat{\beta} = 58.6$~ms, $SE = 54.2$, $t = 1.08$, $p = .28$).

Critically, the effect was strongly modulated by language (H3):
the Punjabi $\times$ passive interaction was significant
($\hat{\beta} = -151.9$~ms, $SE = 76.7$, $t = -1.98$, $p = .048$),
indicating that Punjabi passive sentences did not carry the same prolonged
fixation cost as English passive sentences
(Punjabi $\Delta$ passive: $M = +22.5$~ms vs.\ English $\Delta$ passive:
$M = +94.9$~ms).
In Punjabi, the three sentence types showed similar mean fixation durations
(active: $M = +26.9$~ms; passive: $M = +22.5$~ms; passive-dropped-agent:
$M = +26.9$~ms), suggesting that Punjabi sentence structure did not
differentially burden fixation duration in the way English passive did.
Descriptive means are given in Table~\ref{tab:trial_means}.

\subsubsection{Fixation Count}

Passive sentences produced significantly fewer fixations per trial than
active sentences (English $\Delta M = -2.12$ vs.\ $-0.19$), providing a
complementary measure of processing cost: listeners made fewer, longer
individual fixations under non-agent-first structure.
The GAMM yielded the largest parametric effect in the trial-level battery:
$\hat{\beta} = -1.93$, $SE = 0.48$, $t = -4.04$, $p < .0001$.
Passive-with-dropped-agent was also significantly reduced relative to active:
$\hat{\beta} = -1.01$, $SE = 0.48$, $t = -2.11$, $p = .035$.

The Language $\times$ Passive interaction replicated the mean-duration finding:
Punjabi passive showed a smaller fixation-count reduction than English passive
($\hat{\beta} = +1.58$, $SE = 0.68$, $t = 2.34$, $p = .020$), again indicating
that the non-agent-first cost in fixation behaviour was attenuated in Punjabi.
No significant interaction was found for passive-with-dropped-agent
($p = .48$), suggesting the agent-drop penalty was similar in both languages.

The dissociation between fixation duration (longer) and fixation count
(fewer) in passive sentences is a well-established signature of
increased per-fixation processing load: rather than making more fixations,
listeners devoted longer processing time to each fixation, consistent with
effortful real-time parsing.

\subsubsection{Saccade Rate}

Passive sentences produced significantly fewer saccades per second than
active sentences ($\Delta$ passive: $M = -1.25$~sac/s vs.\ $\Delta$ active:
$M = -1.00$~sac/s; $\hat{\beta} = -0.25$~sac/s, $SE = 0.07$, $t = -3.64$,
$p < .001$), converging with the fixation-count reduction.
Note that the intercept (active condition baseline) was itself strongly
negative relative to free-viewing ($\hat{\beta} = -1.00$, $SE = 0.07$,
$t = -14.67$, $p < .0001$), confirming that sentence-listening globally
reduces scanning behaviour; the passive penalty is a further suppression
above this global effect.
The passive-with-dropped-agent condition was numerically lower than active
but not significantly so ($p = .27$), failing to provide clear evidence
for H2 on this measure.
The Language $\times$ Passive interaction was significant
($\hat{\beta} = +0.21$~sac/s, $SE = 0.10$, $t = 2.16$, $p = .031$):
the saccade-rate reduction under passive was partially recovered in Punjabi,
mirroring the fixation-duration and fixation-count interactions.

\subsubsection{Saccade Duration}

Saccadic movements during passive sentence processing were significantly
longer in duration than during active processing
($\hat{\beta} = +22.95$~ms, $SE = 10.59$, $t = 2.17$, $p = .030$;
English $\Delta M$ passive: $+27.3$~ms vs.\ active: $+4.4$~ms),
indicating slower oculomotor transitions under non-agent-first structure.
No significant language modulation was found for this measure ($p = .31$).

\subsubsection{Total Fixation Duration and Dwell Time}

Neither total fixation duration nor normalised interest-area dwell time
showed significant sentence-type effects in the GAMMs (all parametric $p > .13$).
However, a significant Language $\times$ Passive interaction was found for
total interest-area dwell time ($\hat{\beta} = -279.5$~ms,
$SE = 126.1$, $t = -2.22$, $p = .027$): Punjabi passive sentences directed
substantially less total dwell time to the interest areas than English passive,
again consistent with the broader pattern of Punjabi mitigating the passive
processing cost.

\subsubsection{Cognitive Load Composite Index}

A composite cognitive load index was computed as the mean of five
standardised $\Delta$ indicators (mean fixation duration, total fixation
duration, fixation count, interest-area dwell time, and negated saccade
amplitude; see \S\ref{sec:processing}).
The composite did not yield significant parametric effects in the GAMM
($p > .35$ for all terms), reflecting the mixed pattern of significance
across its constituent measures.
The composite is reported as secondary converging evidence only.

\begin{table}[h]
\centering
\caption{Descriptive means ($M$) and standard deviations ($SD$) for
         $\Delta$ trial-level outcomes by language and sentence type.
         All values are baseline-adjusted (sentence-audio minus free-viewing).}
\label{tab:trial_means}
\begin{threeparttable}
\begin{tabular}{llrrrrrr}
\toprule
& & \multicolumn{2}{c}{Active} & \multicolumn{2}{c}{Passive} & \multicolumn{2}{c}{Passive-dropped} \\
\cmidrule(lr){3-4}\cmidrule(lr){5-6}\cmidrule(lr){7-8}
Language & Outcome & $M$ & $SD$ & $M$ & $SD$ & $M$ & $SD$ \\
\midrule
\multirow{4}{*}{English}
  & $\Delta$ Mean fix. dur.\ (ms)     & $-52.9$ & 511.3 & $+94.9$ & 474.9 & $+5.8$  & 539.5 \\
  & $\Delta$ Fixation count           & $-0.19$ &   4.4 & $-2.12$ &   4.6 & $-1.19$ &   4.3 \\
  & $\Delta$ Saccade rate (sac/s)     & $-1.00$ &   0.7 & $-1.25$ &   0.8 & $-1.07$ &   0.6 \\
  & $\Delta$ Saccade duration (ms)    & $+4.4$  &  25.1 & $+27.3$ & 212.3 & $+13.6$ &  34.6 \\[4pt]
\multirow{4}{*}{Punjabi}
  & $\Delta$ Mean fix.\ dur.\ (ms)    & $+26.9$ & 488.9 & $+22.5$ & 546.3 & $+26.9$ & 310.2 \\
  & $\Delta$ Fixation count           & $-0.38$ &   4.1 & $-0.73$ &   4.6 & $-0.92$ &   4.1 \\
  & $\Delta$ Saccade rate (sac/s)     & $-1.00$ &   0.6 & $-1.04$ &   0.7 & $-1.11$ &   0.6 \\
  & $\Delta$ Saccade duration (ms)    & $+9.8$  &  37.2 & $+17.6$ &  37.4 & $+11.4$ &  36.5 \\
\bottomrule
\end{tabular}
\begin{tablenotes}
\footnotesize
\item Note: Passive-dropped = passive-with-dropped-agent.
      $\Delta$ values = sentence-audio minus matched free-viewing baseline.
      Positive values indicate more gaze activity than baseline.
\end{tablenotes}
\end{threeparttable}
\end{table}

\begin{table}[h]
\centering
\caption{GAMM results for trial-level outcomes.
         Reference levels: English, active.
         Random effects: $s(\text{ParticipantID})$, $s(\text{ImageID})$.}
\label{tab:gamm_trial}
\begin{threeparttable}
\begin{tabular}{llrrrr}
\toprule
Outcome & Term & $\hat{\beta}$ & $SE$ & $t$ & $p$ \\
\midrule
\multirow{5}{*}{$\Delta$ Mean fix.\ dur.\ (ms)}
  & Intercept                            & $-52.78$ & 42.42 & $-1.24$ & .214 \\
  & Punjabi                              & $+79.54$ & 54.11 & $+1.47$ & .142 \\
  & Passive                              & $\mathbf{+147.63}$ & 54.30 & $+2.72$ & $\mathbf{.007}$ \\
  & Passive-dropped                      & $+58.58$ & 54.21 & $+1.08$ & .280 \\
  & Punjabi $\times$ Passive             & $\mathbf{-151.87}$ & 76.66 & $-1.98$ & $\mathbf{.048}$ \\[4pt]
\multirow{5}{*}{$\Delta$ Fixation count}
  & Intercept                            & $-0.19$ & 0.39 & $-0.49$ & .622 \\
  & Punjabi                              & $-0.19$ & 0.48 & $-0.40$ & .692 \\
  & Passive                              & $\mathbf{-1.93}$ & 0.48 & $-4.04$ & $\mathbf{<.001}$ \\
  & Passive-dropped                      & $\mathbf{-1.01}$ & 0.48 & $-2.11$ & $\mathbf{.035}$ \\
  & Punjabi $\times$ Passive             & $\mathbf{+1.58}$ & 0.68 & $+2.34$ & $\mathbf{.020}$ \\[4pt]
\multirow{4}{*}{$\Delta$ Saccade rate (sac/s)}
  & Intercept                            & $-1.00$ & 0.07 & $-14.67$ & $<.0001$ \\
  & Passive                              & $\mathbf{-0.25}$ & 0.07 & $-3.64$ & $\mathbf{.0003}$ \\
  & Punjabi $\times$ Passive             & $\mathbf{+0.21}$ & 0.10 & $+2.16$ & $\mathbf{.031}$ \\[4pt]
\multirow{2}{*}{$\Delta$ Saccade duration (ms)}
  & Passive                              & $\mathbf{+22.95}$ & 10.59 & $+2.17$ & $\mathbf{.030}$ \\[4pt]
\multirow{1}{*}{$\Delta$ IA dwell time (ms)}
  & Punjabi $\times$ Passive             & $\mathbf{-279.48}$ & 126.09 & $-2.22$ & $\mathbf{.027}$ \\
\bottomrule
\end{tabular}
\begin{tablenotes}
\footnotesize
\item Note: Only theoretically relevant terms are shown; see supplementary
      materials for full model tables.
      Bold = $p < .05$.
      Passive-dropped = passive-with-dropped-agent.
\end{tablenotes}
\end{threeparttable}
\end{table}

% ============================================================
\subsection{Interest-Area Attention Allocation (H4)}
% ============================================================

\subsubsection{AOI Dwell Time: Object Dominance}

Across all six conditions, the object/patient AOI received more dwell time
than the agent AOI (Table~\ref{tab:aoi_means}).
This advantage for the object AOI was consistent and reached
significance in multiple conditions after Benjamini--Hochberg correction.
In English active sentences, object dwell exceeded agent dwell by
344~ms ($M_{\text{object}} = 2171$~ms, $M_{\text{agent}} = 1827$~ms;
Mann--Whitney $U$, $p_{\text{BH}} = .041$).
The object advantage was larger and more reliable in English passive
($\Delta = 599$~ms; $p_{\text{BH}} = .002$) and English passive-with-dropped-agent
($\Delta = 499$~ms; $p_{\text{BH}} = .002$).
In Punjabi, a significant object advantage emerged for passive-with-dropped-agent
($\Delta = 440$~ms; $p_{\text{BH}} = .002$) but not for active ($\Delta = 123$~ms;
$p_{\text{BH}} = .16$) or passive ($\Delta = 13$~ms; $p_{\text{BH}} = .14$).

The GAMM for AOI dwell time (normalized) did not yield significant parametric
effects for sentence type or AOI role ($p > .12$), but the AOI-label random
smooth was massively significant ($F_{32,34} = 17.2$, $p < .001$), indicating
that individual image content accounted for far more variance in where
participants looked than sentence structure did.
This finding underscores the importance of the AOI-level random effects
in the model.

\subsubsection{AOI Entry Count}

The GAMM for entry count (the number of times gaze entered each AOI)
showed that passive sentences produced significantly fewer AOI entries than
active sentences ($\hat{\beta} = -0.26$, $SE = 0.12$, $t = -2.12$, $p = .034$),
consistent with the trial-level fixation-count reduction.

\subsubsection{Time to First Fixation in AOI}

In English passive-with-dropped-agent sentences, the object AOI was reached
significantly sooner than the agent AOI (object: $M = 6085$~ms vs.\ agent:
$M = 6356$~ms; $p_{\text{BH}} = .016$).
When the agent is entirely absent from the sentence, listeners oriented gaze
first toward the object of the action rather than the site of the (absent)
agent.
No significant agent/object difference in time to first fixation was found
for active or passive sentences, consistent with both AOIs being approximately
equally salient during those conditions.
The GAMM parametric terms for time to first fixation were non-significant
(all $p > .08$), though the marginal trend for passive-with-dropped-agent
($\hat{\beta} = +145$~ms toward the agent, $t = 1.73$, $p = .085$) suggests
a possible numeric delay before agent fixation in that condition.

\begin{table}[h]
\centering
\caption{Mean dwell time (ms), normalised dwell (proportion of trial),
         entry count, and time to first fixation (ms) by language, sentence type,
         and AOI role (agent vs.\ object).}
\label{tab:aoi_means}
\begin{threeparttable}
\begin{tabular}{lllrrrr}
\toprule
Language & Sentence Type & AOI & Dwell (ms) & Norm.\ dwell & Entries & TFF (ms) \\
\midrule
\multirow{6}{*}{English}
  & \multirow{2}{*}{Active}
    & Agent  & 1827 & .42 & 2.53 & 6189 \\
  & & Object & 2171$^{*}$ & .49$^{*}$ & 2.65 & 6097 \\[2pt]
  & \multirow{2}{*}{Passive}
    & Agent  & 1712 & .39 & 2.26 & 6232 \\
  & & Object & 2311$^{**}$ & .52$^{**}$ & 2.36 & 6049 \\[2pt]
  & \multirow{2}{*}{Passive-dropped}
    & Agent  & 1737 & .40 & 2.32 & 6356 \\
  & & Object & 2235$^{**}$ & .51$^{**}$ & 2.40 & 6085$^{*}$ \\
\midrule
\multirow{6}{*}{Punjabi}
  & \multirow{2}{*}{Active}
    & Agent  & 1943 & .45 & 2.69 & 6121 \\
  & & Object & 2066 & .47 & 2.59 & 6164 \\[2pt]
  & \multirow{2}{*}{Passive}
    & Agent  & 1952 & .45 & 2.50 & 6131 \\
  & & Object & 1965 & .46 & 2.57 & 6061 \\[2pt]
  & \multirow{2}{*}{Passive-dropped}
    & Agent  & 1740 & .40 & 2.46 & 6309 \\
  & & Object & 2181$^{**}$ & .50$^{**}$ & 2.54 & 6060 \\
\bottomrule
\end{tabular}
\begin{tablenotes}
\footnotesize
\item Note: Superscripts indicate significant difference between agent and
      object for that condition (Mann--Whitney, BH-corrected):
      $^{*}p < .05$; $^{**}p < .01$.
      TFF = time to first fixation (ms from trial onset).
      Passive-dropped = passive-with-dropped-agent.
\end{tablenotes}
\end{threeparttable}
\end{table}

% ============================================================
\subsection{First Non-Centre AOI Visit (H4)}
% ============================================================

After excluding the centre-of-screen fixation (index~1), the first fixation
landing on a coded AOI was extracted for each trial.
The proportion of trials on which this first meaningful AOI visit was directed
to the agent AOI is shown in Table~\ref{tab:first_look}.
Across all conditions, listeners looked first to the \emph{object} AOI more
often than to the agent AOI (proportions below 0.5 in all but the Punjabi passive condition).

A directional but non-significant trend emerged consistent with H4:
in English, the probability of an initial agent look decreased monotonically
from active (37.3\%) to passive (35.8\%) to passive-with-dropped-agent (31.8\%),
indicating that removing the agent from the sentence reduced initial agent
orientation.
In Punjabi, the passive condition was a striking exception:
participants directed their first AOI fixation to the agent on 49.0\% of
passive trials, approaching chance and substantially exceeding the English
passive rate (35.8\%).
This cross-linguistic dissociation --- Punjabi passive resembles agent-first
attention allocation, English passive does not --- is consistent with
Punjabi morphological cues maintaining agent salience even in passive structures.

A logistic GAMM confirmed that no individual contrast reached statistical
significance (all $|z| < 1.30$, all $p > .19$), likely reflecting
insufficient power for a binary first-look measure at $n = 148$--$150$ trials
per cell.
The predicted probabilities from this model are shown in Table~\ref{tab:first_look}
alongside raw proportions.

\begin{table}[h]
\centering
\caption{Proportion of trials with agent-directed first AOI fixation
         (after excluding fixation index~1).}
\label{tab:first_look}
\begin{threeparttable}
\begin{tabular}{llrrrrr}
\toprule
& & \multicolumn{2}{c}{Observed} & \multicolumn{3}{c}{GAMM-predicted} \\
\cmidrule(lr){3-4}\cmidrule(lr){5-7}
Language & Sentence Type & $P(\text{agent first})$ & $N$ & $\hat{P}$ & 95\% CI \\
\midrule
English & Active                & .373 & 150 & .342 & [.19, .50] \\
English & Passive               & .358 & 148 & .322 & [.17, .47] \\
English & Passive-dropped       & .318 & 148 & .279 & [.14, .42] \\
Punjabi & Active                & .413 & 150 & .398 & [.23, .56] \\
Punjabi & Passive               & .490 & 149 & .499 & [.33, .67] \\
Punjabi & Passive-dropped       & .405 & 148 & .399 & [.23, .56] \\
\bottomrule
\end{tabular}
\begin{tablenotes}
\footnotesize
\item Note: No contrast reached significance in the logistic GAMM ($p > .19$).
      Passive-dropped = passive-with-dropped-agent.
      GAMM predictions exclude random effects for participant and image.
\end{tablenotes}
\end{threeparttable}
\end{table}

% ============================================================
\subsection{Temporal Gaze Dynamics}
% ============================================================

To examine how attention allocation evolved across the trial, the proportion
of agent-directed fixations was computed at each ordinal fixation position
(index 2 through 17) for each condition.
Bins containing fewer than 20 fixations were excluded.
The resulting curves (Figure~\ref{fig:temporal}) and their smoothed GAMM fits
reveal several patterns consistent with the agent-first principle.

In English, active sentences showed a rising trajectory of agent-directed looks:
beginning at 36.5\% at position~2 and reaching approximately 55\% by positions
11--12 before declining.
This rise indicates that as the active sentence unfolds, listeners increasingly
orient toward the agent, consistent with active deployment of an agent-first
parsing strategy.
English passive sentences, by contrast, showed a flatter and lower trajectory,
remaining between 36\% and 49\% throughout the trial, without the sustained
rise observed for active.
English passive-with-dropped-agent started lowest at position~2 (32.1\%) and
remained suppressed through positions~2--4 before recovering, suggesting
initial suppression of agent attention when no overt agent was mentioned.

The Punjabi patterns diverged from English in theoretically informative ways.
Punjabi passive sentences showed the highest agent-look proportions of any
condition at early fixation positions (49.6\% at index~2), exceeding Punjabi
active (39.7\%) at the same position.
This early Punjabi passive advantage is consistent with Punjabi morphological
case marking rapidly communicating agent identity even in passive structures,
enabling listeners to orient gaze to the agent promptly.
By late positions (indices 12--13), Punjabi active overtook Punjabi passive,
suggesting that the agent-first advantage for active structures emerges
more gradually in Punjabi than in English.

These temporal dynamics provide a finer-grained picture of the agent-first
principle than aggregate fixation summaries alone:
English listeners rapidly commit to an active agent-first parsing strategy,
whereas Punjabi listeners show less differentiation early but equivalent or
greater agent orientation when morphological cues are available (as in Punjabi passive).

% ============================================================
\subsection{Summary of Hypothesis Tests}
% ============================================================

\begin{table}[h]
\centering
\caption{Summary of hypothesis tests.}
\label{tab:hypotheses}
\begin{threeparttable}
\begin{tabular}{p{0.5cm} p{5.5cm} p{5cm} p{2cm}}
\toprule
 & Hypothesis & Key evidence & Verdict \\
\midrule
H1 & Passive sentences impose greater gaze-based processing cost than active &
  $\Delta$ mean fix.\ dur.\ $+148$~ms in passive (English; $p=.007$);
  $\Delta$ fix.\ count $-1.93$ ($p<.001$);
  $\Delta$ saccade rate $-0.25$ sac/s ($p<.001$) &
  \textbf{Supported} \\[6pt]
H2 & Passive-dropped-agent imposes at least as great a cost as passive-with-by-phrase &
  Fix.\ count reduction $-1.01$ vs.\ $-1.93$ (both significant);
  mean fix.\ dur.\ non-significant for passive-dropped;
  no consistent ordering across measures &
  \textbf{Partial} \\[6pt]
H3 & Cross-linguistic modulation: Punjabi reduces the passive cost &
  Three significant Language $\times$ Passive interactions
  (mean fix.\ dur., fix.\ count, saccade rate, IA dwell time);
  Punjabi passive approximates Punjabi active on all measures &
  \textbf{Supported} \\[6pt]
H4 & Agent-first gaze bias: active promotes earlier/more agent looking &
  English monotonic decrease in $P(\text{agent first look})$: active $>$ passive $>$ passive-dropped;
  temporal curves show active agent-look rise absent for passive;
  Punjabi dissociation is consistent with morphology-driven agent salience &
  \textbf{Directionally supported; underpowered} \\
\bottomrule
\end{tabular}
\end{threeparttable}
\end{table}

% ============================================================
\section*{Supplementary Notes}
% ============================================================

\paragraph{Random effects.}
Participant random effects were significant in all trial-level GAMMs
($p < .001$ in all models), confirming large individual differences in
baseline gaze behaviour.
Image random effects were significant in the saccade-rate model
($p = .002$) and all interest-area models ($p < .001$), but not in
fixation duration or fixation count models, indicating that stimulus
images varied more in how they modulated scanning than in fixation duration.
The AOI-label random smooth captured enormous variance in all four
interest-area models ($F > 7$, $p < .001$), indicating that individual
AOI regions (specific image regions) were far stronger predictors of
dwell time than sentence structure.

\paragraph{Model fit.}
Deviance explained ranged from 2.6\% (normalised dwell time) to 33.6\%
(AOI entry count), with saccade rate (27.5\%) and interest-area entry
count (33.6\%) showing the best overall fit.
Trial-level fixation models explained 10--15\% of deviance, reflecting
the high between-trial variability in gaze behaviour.

\paragraph{Centre-bias correction.}
Excluding fixation index~1 from trial-level aggregation produced
modestly cleaner effect estimates: the passive-dropped-agent condition
in English moved from $+63$~ms (with index~1 included) to $+6$~ms
after exclusion, suggesting that the stereotyped first fixation had
been partially obscuring condition differences for the dropped-agent condition
specifically.
The passive vs.\ active contrast for mean fixation duration
(+95~ms after correction) was robust to this correction and remained
significant.

\bigskip

\bibliographystyle{apa}
\bibliography{references}

\end{document}
